\chapter{Real and Complex Numbers: Foundations of Analysis}

\section{Real Number Systems and the Real Line}

\subsection{Natural Numbers, Integers, and Rational Numbers}

The real number system, $\mathbb{R}$, can be approached through a sequence of progressively larger number systems, each containing the preceding one.
The natural numbers form the starting point:
\[
\mathbb{N} = \{0, 1, 2, 3, \dots\}.
\]
Including their negatives gives the integers:
\[
\mathbb{Z} = \{0, 1, -1, 2, -2, \dots\}.
\]
Allowing quotients of integers with nonzero denominators then gives the rational numbers, $\mathbb{Q}$:
\[
\mathbb{Q} = \left\{ \pm \frac{p}{q} : p \in \mathbb{N}, \ q \in \mathbb{N}, \ q \neq 0 \right\}.
\]

\subsection{Irrational, Algebraic, and Transcendental Numbers}

The rational numbers do not exhaust the real line. Real numbers that cannot be expressed as ratios of integers are called irrational numbers.
The following theorem provides a fundamental example.

\begin{theorem}[Irrationality of the Square Root of Two]
The square root of $2$ is irrational.
\end{theorem}

\begin{proof}
Suppose, for a contradiction, that $\sqrt{2}$ is rational. Then we can write $\sqrt{2} = p/q$, for integers $p, q$ with $q \neq 0$. Assume that $p/q$ is in lowest terms; in particular, $p$ and $q$ cannot both be even.
Squaring both sides gives $2 q^2 = p^2$. Thus $p^2$ is even. Since the square of an odd integer is odd, $p$ must also be even.
Writing $p = 2k$ for an integer $k$ and substituting gives
\[
\begin{aligned}
2q^2 &= 4k^2  \\
\Rrightarrow\quad  q^2 &= 2k^2.
\end{aligned}
\]

Hence $q^2$ is even, so $q$ is even as well. Both $p$ and $q$ are therefore even, contradicting the assumption that $p/q$ is in lowest terms.
This proves the theorem.
\end{proof}

The real numbers comprise both rational and irrational numbers. Familiar irrational numbers include $\sqrt{2}$, $\pi$, Euler's number $e$,
and the golden ratio $\varphi$. For two lengths $a$ and $b$ in that ratio, the defining proportion is
\[
\varphi = \frac{a+b}{a} = \frac{a}{b}.
\]

The numbers $\sqrt{2}$ and $\varphi$ also belong to the class of algebraic numbers. A real number $x$ is algebraic if it is a root of a
nonzero polynomial with integer coefficients $a_k \in \mathbb{Z}$, for $k = 0, \dots, n$:
\[
\sum_{k=0}^{n} a_k x^k = 0.
\]

Both $\sqrt{2}$ and $\varphi$ satisfy this definition. By contrast, $\pi$ is transcendental: it is not algebraic.
The historical results on $\pi$ established between 1794 and 1882 illustrate the distinction between irrationality and transcendence.
In particular, irrational algebraic numbers form a proper subclass of the irrational numbers.

\subsection{The Real Line and Interval Notation}

The real line represents $\mathbb{R}$ geometrically: each real number corresponds to a point, and each point corresponds to a real number.
Its absence of gaps reflects the completeness property discussed below. Intervals are particularly important subsets of the real line;
the following notation will be used throughout the text.
\[
\begin{aligned}
\relax
[a,b] &= \{x \in \mathbb{R} : a \leq x \leq b\}, \\
(a,b) &= \{x \in \mathbb{R} : a < x < b\}, \\
[a,b) &= \{x \in \mathbb{R} : a \leq x < b\}, \\
(a, \infty) &= \{x \in \mathbb{R} : x > a\}, \\
(-\infty, b] &= \{x \in \mathbb{R} : x \leq b\}.
\end{aligned}
\]

\begin{remark}
The symbol $\infty$ is not a real number. In interval notation, it indicates that an interval is unbounded in the corresponding direction.
\end{remark}

Figure~\ref{fig:real-intervals} distinguishes included and excluded endpoints and shows how rays represent unbounded intervals.

\begin{figure}[!htbp]
\centering
\begin{tikzpicture}[x=1.3cm,y=1cm,>=Stealth]
  \foreach \rowheight in {0,-1.2,-2.4,-3.6,-4.8} {
    \draw[->,black!35] (0,\rowheight) -- (6,\rowheight);
  }
  \node[left] at (-0.3,0) {$[a,b]$};
  \node[left] at (-0.3,-1.2) {$(a,b)$};
  \node[left] at (-0.3,-2.4) {$[a,b)$};
  \node[left] at (-0.3,-3.6) {$(a,\infty)$};
  \node[left] at (-0.3,-4.8) {$(-\infty,b]$};
  \foreach \rowheight in {0,-1.2,-2.4} {
    \draw[figblue,very thick] (1,\rowheight) -- (4.5,\rowheight);
    \node[below=5pt] at (1,\rowheight) {$a$};
    \node[below=5pt] at (4.5,\rowheight) {$b$};
  }
  \fill[figblue] (1,0) circle (2.5pt) (4.5,0) circle (2.5pt);
  \draw[figblue,thick,fill=white] (1,-1.2) circle (2.5pt) (4.5,-1.2) circle (2.5pt);
  \fill[figblue] (1,-2.4) circle (2.5pt);
  \draw[figblue,thick,fill=white] (4.5,-2.4) circle (2.5pt);
  \draw[figgreen,very thick,->] (1,-3.6) -- (5.8,-3.6);
  \draw[figgreen,thick,fill=white] (1,-3.6) circle (2.5pt);
  \node[below=5pt] at (1,-3.6) {$a$};
  \draw[figgreen,very thick,->] (4.5,-4.8) -- (0.2,-4.8);
  \fill[figgreen] (4.5,-4.8) circle (2.5pt);
  \node[below=5pt] at (4.5,-4.8) {$b$};
\end{tikzpicture}
\caption[Intervals on the real line]{Closed, open, and half-open intervals, followed by two unbounded rays. Filled endpoints belong to the set; hollow endpoints do not. Colored arrows indicate continuation without bound, rather than an endpoint at infinity.}
\label{fig:real-intervals}
\end{figure}

\section{Field Structure, Order, and Completeness}

The real number system can also be characterized axiomatically. Three groups of properties describe its algebraic operations, order, and completeness.

\begin{definition}[The Real Number System]
$\mathbb{R}$ is a complete ordered field.
\end{definition}

\subsection{Axioms of Addition and Multiplication}

We begin with the notion of a field. Addition and multiplication are binary operations on $\mathbb{R}$ satisfying the following laws.
\begin{itemize}
\item \textit{Associative laws}: for all real numbers $a$, $b$, and $c$, both identities hold:

\[
\begin{aligned}
a + (b+c) &= (a+b) + c \\
a(bc) &= (ab)c.
\end{aligned}
\]

\item \textit{Commutative laws}: for all real numbers $a$, $b$, and $c$, both identities hold:

\[
\begin{aligned}
a + b &= b + a \\
ab &= ba.
\end{aligned}
\]

\item \textit{Identity laws}: for every real number $a$, the additive and multiplicative identities satisfy

\[
\begin{aligned}
a + 0 &= a \\
a \cdot 1 &= a.
\end{aligned}
\]

\item \textit{Inverse laws}: for every real number $a$, there is a unique real number $-a$ such that $a + (-a) = 0$, and, if $a \neq 0$, a unique real number $a^{-1}$ such that $a \, a^{-1} = 1$.

\item \textit{Distributive law}: for all real numbers $a$, $b$, and $c$, multiplication distributes over addition:

\[
a(b+c) = ab + ac.
\]

\end{itemize}

\subsection{Order and Compatibility with Arithmetic}

The field $\mathbb{R}$ is ordered: its strict order relation $<$ compares any two distinct real numbers and is compatible with addition and multiplication.
The order satisfies the following properties.
\begin{itemize}
\item \textit{Transitivity}: if $x < y$ and $y < z$ for real numbers $x, y, z$, then $x < z$.
\item \textit{Trichotomy}: for any two real numbers $x, y$, exactly one of the following relations holds:

\[
x < y, \quad x = y, \quad y < x.
\]

\end{itemize}

Compatibility with the field operations means that:
\begin{itemize}
\item \textit{Addition}: if $x < y$, then $x + z < y + z$ for every real number $z$.
\item \textit{Multiplication}: if $x < y$ and $z > 0$, then $xz < yz$.
\end{itemize}

\subsection{Completeness and the Least Upper Bound Property}

Completeness is expressed by the \emph{Axiom of Continuity}. Among its equivalent formulations, we use the least upper bound property.
To state it, we first introduce the notion of an upper bound. Let $A \subset \mathbb{R}$ be nonempty.

\begin{definition}[Upper Bounds]
A real number $b$ is an upper bound of $A$ if $x \leq b$ for every $x \in A$. A set with an upper bound is said to be bounded above.
\end{definition}

The Axiom of Continuity can now be stated precisely.

\begin{theorem}[Axiom of Continuity: Existence of the Least Upper Bound]
Every nonempty subset $A \subset \mathbb{R}$ that admits an upper bound has a least upper bound in $\mathbb{R}$.
\end{theorem}

This property distinguishes the real numbers from the rational numbers: it does not hold in $\mathbb{Q}$. For example, the set $A = \{ x \in \mathbb{Q} : x \geq 0, \ x^2 < 2 \} \subset \mathbb{Q}$ is bounded above in $\mathbb{Q}$, for example by $2$, but has no least upper bound in $\mathbb{Q}$. Such a bound would have to be $\sqrt{2}$,
which is irrational. In $\mathbb{R}$, the least upper bound guaranteed by the Axiom of Continuity has a special name and notation.

\subsection{Suprema, Maxima, and Infima}

\begin{definition}[The Supremum]
Let $A \subset \mathbb{R}$ be nonempty and bounded above. Its supremum, denoted $b = \sup A$, is its least upper bound.
Equivalently, $b = \sup A$ if and only if $x \leq b$ for every $x \in A$ and, for every $\varepsilon > 0$, there exists
$x_\varepsilon \in A$ such that $x_\varepsilon > b - \varepsilon$. The first condition makes $b$ an upper bound; the second ensures that no smaller number is an upper bound.
\end{definition}

Two distinctions are useful.
\begin{enumerate}
\item \textit{Maximum}: a real number $b$ is the maximum of $A$, denoted $b = \max A$, if $x \leq b$ for every $x \in A$ and $b \in A$.
A supremum need not belong to the set. For example, $A = [0,1)$ has supremum $1$ but no maximum, since $1 \notin A$.
Whenever a maximum exists, it equals the supremum: $\max A = \sup A$.
\item \textit{Unbounded sets}: for a set $A \subset \mathbb{R}$ that is not bounded above, the convention is to write $+\infty := \sup A$.
This extends the notation beyond sets with a finite supremum.
\end{enumerate}

Reversing the inequalities gives the equivalent greatest lower bound formulation of the Axiom of Continuity.

\begin{theorem}[Axiom of Continuity: Existence of the Greatest Lower Bound]
Every nonempty subset $A \subset \mathbb{R}$ that admits a lower bound has a greatest lower bound, called the infimum of $A$, in $\mathbb{R}$.
\end{theorem}

In summary, $\mathbb{R}$ is a field because its arithmetic operations satisfy the field axioms. It is ordered because these operations are compatible
with its order. In particular, addition preserves order for every $z$:
\[
x < y \quad \text{implies} \quad x+z < y+z.
\]
Likewise, the product of nonnegative numbers is nonnegative:
\[
x \geq 0 \quad \text{together with} \quad y \geq 0 \quad \text{implies} \quad xy \geq 0.
\]

Finally, $\mathbb{R}$ is complete because it satisfies the Axiom of Continuity.
Together, these properties characterize the real number system.

Figure~\ref{fig:atlas-supremum} contrasts a supremum that is not attained with a maximum that belongs to the set.

\begin{figure}[!htbp]
\centering
\begin{tikzpicture}[>=Stealth]
\draw[->,black!45] (-0.3,0) -- (6.5,0);
\draw[figblue,very thick] (0,0) -- (5,0);
\fill[figblue] (0,0) circle (2pt);
\draw[figblue,thick,fill=white] (5,0) circle (2.5pt);
\node[below] at (0,0) {$0$};
\node[below] at (5,0) {$1$};
\node[left] at (-0.5,0) {$[0,1)$};
\draw[dashed,figgreen] (3.5,-0.2) -- (3.5,0.7);
\node[above,figgreen] at (3.5,0.7) {$1-\varepsilon$};
\fill[figred] (4.25,0) circle (2.5pt);
\draw[figred] (4.25,0.1) -- (4.25,1.3);
\node[above,figred] at (4.25,1.3) {$1-\varepsilon/2$};
\draw[->,black!45] (-0.3,-1.7) -- (6.5,-1.7);
\draw[figblue,very thick] (0,-1.7) -- (5,-1.7);
\fill[figblue] (0,-1.7) circle (2pt) (5,-1.7) circle (2.5pt);
\node[below] at (0,-1.7) {$0$};
\node[below] at (5,-1.7) {$1$};
\node[left] at (-0.5,-1.7) {$[0,1]$};
\end{tikzpicture}
\caption{The sets $[0,1)$ (top) and $[0,1]$ (bottom) both have supremum $1$. Only the second contains that endpoint and therefore has a maximum. For the illustrated choice $0<\varepsilon<1$, the highlighted point $1-\varepsilon/2$ lies within $\varepsilon$ of the supremum in the first set.}
\label{fig:atlas-supremum}
\end{figure}

\section{Decimal Expansions and Geometric Series}

\subsection{Decimal Representations as Infinite Series}

Real numbers can also be represented in decimal notation. Every real number $x$ can be written in the form $x = \pm b_0.b_1 b_2 \ldots$, where $b_0 \in \mathbb{N}$ gives the digits before the decimal point, and the digits $b_i \in \{0, 1, \dots, 9\}$ form the fractional part.
A decimal expansion may terminate, repeat periodically, or continue indefinitely without repeating. The first two examples below are repeating decimals; the third is the nonrepeating expansion of $\pi$, which has been computed to more than thirty-one trillion digits:
\[
\begin{aligned}
0.33333\ldots &= 0.\overline{3}, \\
0.142857142857\ldots &= 0.\overline{142857}, \\
\pi &= 3.1415926535897\ldots.
\end{aligned}
\]
Thus, real numbers admit finite, repeating, or nonrepeating decimal representations. To give this notation a precise meaning, we interpret a decimal expansion as the value of an infinite series, defined as the limit of its partial sums:
\[
\begin{aligned}
x &= b_0.b_1 b_2 \ldots \\
&= b_0 + \frac{b_1}{10} + \frac{b_2}{100} + \cdots + \frac{b_n}{10^n} + \cdots \\
&= \sum_{k=0}^{\infty} \frac{b_k}{10^k} \\
&:= \lim_{n \to \infty} \sum_{k=0}^{n} \frac{b_k}{10^k}.
\end{aligned}
\]

\subsection{Finite Geometric Sums}

The treatment of repeating decimals relies on the formula for the sum of a finite geometric progression.

\begin{theorem}[Geometric Sum Formula]
Given any $q \neq 1$,
\[
\sum_{k=0}^{n-1} q^k = \frac{1 - q^n}{1-q}.
\]

\end{theorem}

\begin{proof}
Multiplying both sides of the identity \[
\sum_{k=0}^{n-1} q^k = q^0 + q^1 + q^2 + \cdots + q^{n-1}
\] by $1-q$, we obtain
\[
\begin{aligned}
(1-q) \sum_{k=0}^{n-1} q^k &= (1-q)(q^0 + q^1 + q^2 + \cdots + q^{n-1}) \\
&= (q^0 + q^1 + \cdots + q^{n-1}) - (q^1 + q^2 + \cdots + q^n).
\end{aligned}
\]
All of the intermediate terms cancel in pairs, leaving only
\[
q^0 - q^n = 1 - q^n,
\]
which proves the formula.
\end{proof}

\subsection{Infinite Geometric Series}

When $|q| < 1$, taking the limit $n \to \infty$ in the finite sum gives the corresponding infinite geometric series.

\begin{theorem}[Geometric Series Formula]
Given any $q \in \mathbb{R}$ with $|q| < 1$,
\[
\sum_{k=0}^{\infty} q^k = \frac{1}{1-q}.
\]

\end{theorem}

\begin{proof}
By definition,
\[
\sum_{k=0}^{\infty} q^k = \lim_{n \to \infty} \sum_{k=0}^{n-1} q^k.
\]
By the Geometric Sum Formula,
\[
\sum_{k=0}^{n-1} q^k = (1 - q^n)/(1-q).
\]
Since $|q| < 1$, the geometric sequence satisfies
\[
\lim_{n \to \infty} q^n = 0,
\]
and the result follows by taking the limit in the finite-sum formula.
\end{proof}

Figure~\ref{fig:atlas-geometric-area} gives an area interpretation of a convergent geometric series.

\begin{figure}[!htbp]
\centering
\begin{tikzpicture}[>=Stealth]
\fill[figblue!28] (0,0) rectangle (2,4);
\fill[figgreen!28] (2,2) rectangle (4,4);
\fill[figred!25] (3,0) rectangle (4,2);
\fill[figblue!50] (2,0) rectangle (3,1);
\draw[thick] (0,0) rectangle (4,4);
\draw (2,0) -- (2,4) (2,2) -- (4,2) (3,0) -- (3,2) (2,1) -- (3,1);
\node at (1,2) {$1/2$};
\node at (3,3) {$1/4$};
\node at (3.5,1) {$1/8$};
\node at (2.5,0.5) {$1/16$};
\node at (2.5,1.5) {$\cdots$};
\draw[<->] (0,-0.4) -- (4,-0.4) node[midway,below] {$1$};
\end{tikzpicture}
\caption{A unit square is successively partitioned into regions of area $1/2$, $1/4$, $1/8$, and $1/16$. The uncolored remainder has area $1/16$ after four steps; continuing the construction makes the remainder tend to zero and illustrates the sum of the geometric series with ratio $1/2$.}
\label{fig:atlas-geometric-area}
\end{figure}

\subsection{Why Repeating Nines Equal One}

The Geometric Series Formula resolves a familiar question about decimal notation: whether $0.\overline{9}$ is equal to $1$.

\begin{theorem}[Repeating Nines: $0.999\ldots=1$]
The repeating decimal $0.\overline{9}=0.999\ldots$ is equal to $1$; that
is,
\[
0.\overline{9}=1.
\]
\end{theorem}

\begin{proof}
By the definition of decimal notation given above, we may write
\[
\begin{aligned}
0.\overline{9} &= 0 + \frac{9}{10} + \frac{9}{100} + \cdots + \frac{9}{10^n} + \cdots \\
&= \frac{9}{10} \sum_{k=0}^{\infty} \frac{1}{10^k}.
\end{aligned}
\]
Applying the Geometric Series Formula with $q = 1/10$, we obtain
\[
\begin{aligned}
\sum_{k=0}^{\infty} (1/10)^k &= 1/(1 - 1/10) \\
&= 10/9,
\end{aligned}
\]
and hence
\[
\begin{aligned}
0.\overline{9} &= \frac{9}{10} \cdot \frac{10}{9} \\
&= 1.
\end{aligned}
\]

\end{proof}

\section{Rational Approximation and Uncountability}

\subsection{Density of the Rational Numbers}

Decimal expansions show that every real number can be approximated arbitrarily well by rational numbers with finite decimal representations.

\begin{theorem}[Density of $\mathbb{Q}$ in $\mathbb{R}$]
Every real number can be approximated to any desired degree of accuracy by rational numbers with finite decimal representations.
\end{theorem}

\begin{proof}
Let $x\in\mathbb{R}$. For each positive integer $n$, define
\[
r_n=\frac{\lfloor 10^n x\rfloor}{10^n}.
\]
The numerator is an integer, so $r_n$ is rational and has a finite
decimal representation. From the defining property of the floor function,
\[
\lfloor 10^n x\rfloor\leq 10^n x<\lfloor 10^n x\rfloor+1.
\]
Dividing by $10^n$ gives
\[
0\leq x-r_n<10^{-n}.
\]
This argument includes $x=0$, negative real numbers, and terminating
decimals; in the terminating case equality may occur on the left. Since
$10^{-n}\to0$, the sequence $r_n$ converges to $x$.
\end{proof}

The same series methods show that the repeating decimal $8.4\overline{35}$ equals
\[
(8435 - 84)/990.
\]

More generally, every repeating decimal represents a rational number, as can be shown by extending the calculation used for $0.\overline{9}$.

\subsection{Countability and Cantor's Diagonal Argument}

Although both $\mathbb{Q}$ and $\mathbb{R}$ are infinite, they differ in size, or \emph{cardinality}.
The rational numbers are countable: there is a bijection, a function that is both one-to-one and onto, from $\mathbb{N}$ to $\mathbb{Q}$.
Consequently, the rational numbers can be listed in an infinite sequence.

The real numbers, by contrast, are uncountable. No one-to-one function from $\mathbb{R}$ to $\mathbb{N}$ exists, so no sequence can list all real numbers.
Even the interval $[0,1]$ is uncountable.

\begin{theorem}[Cantor's Uncountability Theorem]
The set $[0,1]$ is uncountable.
\end{theorem}

\textit{Cantor's diagonal argument}, introduced in 1891, explains the underlying idea.
Consider the set $T$ of infinite binary sequences, whose entries are either zero or one, and suppose that all its elements have been listed.
Construct a new sequence by choosing its $n$-th digit to differ from the $n$-th digit of the $n$-th listed sequence, for every $n$.
The new sequence differs from every entry in the list, contradicting the assumed enumeration.
The full details connecting this construction to the interval are beyond the scope of this discussion.

\section{Complex Numbers and Polynomial Equations}

\subsection{Cartesian Representation and Arithmetic}

Complex numbers extend the real number system to allow solutions of polynomial equations such as $x^2 + 1 = 0$, which has no solution in $\mathbb{R}$.

\begin{definition}[Complex Numbers in Cartesian Form]
A complex number is an expression of the form $a + bi$, where $a$ and $b$ are real numbers, and $i$ is a symbol satisfying $i^2 = -1$. The complex
number $a+bi$ can also be represented by the ordered pair $(a,b)$ and plotted as a point in the \emph{Argand plane}.
In particular, the complex number $i = 0 + 1 \cdot i$ is identified with the point $(0,1)$.
\end{definition}

For $a+bi$, the number $a$ is the real part and $b$ is the imaginary part. For example, $4 - 3i$ has real part $4$ and imaginary part $-3$.
Two complex numbers $a+bi$ and $c+di$ are equal exactly when $a = c$ and $b = d$.
In the Argand plane, the horizontal axis is the real axis and the vertical axis is the imaginary axis. Addition and subtraction are performed separately on the real and imaginary parts:
\[
\begin{aligned}
(a+bi) + (c+di) &= (a+c) + (b+d)i, \\
(a+bi) - (c+di) &= (a-c) + (b-d)i.
\end{aligned}
\]
For instance,
\[
(1-i) + (4+7i) = 5 + 6i.
\]

The product of two complex numbers is defined in such a way that the usual commutative and distributive laws continue to hold, together with the
defining relation $i^2 = -1$; expanding the product $(a+bi)(c+di)$ using distributivity and substituting $i^2 = -1$ yields
\[
(a+bi)(c+di) = (ac - bd) + (ad+bc) i.
\]
As an illustration,
\[
\begin{aligned}
(-1+3i)(2-5i) &= (-1)(2-5i) + 3i(2-5i) \\
&= -2 + 5i + 6i - 15 i^2 \\
&= 13 + 11i.
\end{aligned}
\]
Here the last equality uses $i^2 = -1$.

\subsection{Conjugation, Modulus, and Division}

Division uses the complex conjugate to make the denominator real, much as rationalization removes a square root from a denominator.

\begin{definition}[Complex Conjugation]
For a complex number $z = a+bi$, its complex conjugate is defined to be $\overline{z} = a - bi$.
\end{definition}

To divide by a nonzero complex number, multiply the numerator and denominator by the conjugate of the denominator. For example, to express $(-1+3i)/(2+5i)$ in the form $a+bi$, use the conjugate $2-5i$. The numerator becomes $(-1+3i)(2-5i)$, which was evaluated above. Dividing by the resulting real denominator gives $13/29 + 11i/29$. Geometrically, the complex conjugate $\overline{z}$ is the reflection of $z$ in the real axis. The complex conjugate obeys several elementary
algebraic properties, namely
\[
\begin{aligned}
\overline{z+w} &= \overline{z} + \overline{w} \\
\overline{zw} &= \overline{z}\,\overline{w}.
\end{aligned}
\]
Both identities follow directly from the definitions of the sum and product of complex numbers. The modulus, or absolute value, $|z|$ of a complex
number $z = a+bi$ is defined as its distance from the origin in the Argand plane, that is,
\[
|z| = \sqrt{a^2+b^2}.
\]
It is a direct consequence of the definitions that
\[
\begin{aligned}
z \overline{z} &= (a+bi)(a-bi) \\
&= a^2 + abi - abi - b^2 i^2 \\
&= a^2 + b^2.
\end{aligned}
\]
Consequently, $z\overline{z} = |z|^2$, which explains the division rule. For a nonzero divisor, it can be written as
\[
\begin{aligned}
z/w &= z\overline{w}/(w\overline{w}) \\
&= z\overline{w}/|w|^2.
\end{aligned}
\]

Figure~\ref{fig:complex-conjugation} illustrates why conjugation preserves the modulus while reversing the sign of the imaginary part.

\begin{figure}[!htbp]
\centering
\begin{tikzpicture}[scale=1.2,>=Stealth]
  \draw[->,black!65] (-0.6,0) -- (4.7,0) node[right] {$\mathrm{Re}$};
  \draw[->,black!65] (0,-2.1) -- (0,2.1) node[above] {$\mathrm{Im}$};
  \draw[gridline,dashed] (0,1.6) -- (3,1.6) -- (3,-1.6) -- (0,-1.6);
  \draw[vecA] (0,0) -- (3,1.6);
  \draw[vecB] (0,0) -- (3,-1.6);
  \fill[figblue] (3,1.6) circle (2pt);
  \fill[figred] (3,-1.6) circle (2pt);
  \node[above right,figblue] at (3,1.6) {$z=a+bi$};
  \node[below right,figred] at (3,-1.6) {$\overline{z}=a-bi$};
  \node[below left] at (0,0) {$0$};
  \node[below right] at (3,0) {$a$};
  \node[left] at (0,1.6) {$b$};
  \node[left] at (0,-1.6) {$-b$};
\end{tikzpicture}
\caption[Complex conjugation as reflection]{A complex number and its conjugate are mirror images across the real axis. They have the same real component and opposite imaginary components; their equal distances from the origin express the identity $|\overline{z}|=|z|$.}
\label{fig:complex-conjugation}
\end{figure}

\subsection{Quadratic Equations and Complex Roots}

Since $i^2 = -1$, one may regard $i$ as a square root of $-1$; however, one also has
\[
\begin{aligned}
(-i)^2 &= i^2 \\
&= -1.
\end{aligned}
\]
Consequently, $-i$ is likewise a square root of $-1$. By convention, $i$ is designated the principal square root of $-1$, and one writes $\sqrt{-1} = i$;
more generally, for any positive number $c$, one writes $\sqrt{-c} = \sqrt{c}\, i$.

With this convention, the usual derivation of the quadratic formula for an equation with a nonzero leading coefficient, $ax^2 + bx + c = 0$, remains valid even when the discriminant $b^2 - 4ac$ is negative. The roots are
\[
x = \frac{-b \pm \sqrt{b^2 - 4ac}}{2a}.
\]

\begin{example}
For $x^2+x+1=0$ we have $a=b=c=1$, hence
\[
\Delta=b^2-4ac=1-4=-3.
\]
Using $\sqrt{-3}=i\sqrt{3}$ in the quadratic formula gives
\[
x=\frac{-1\pm\sqrt{-3}}{2}
=\frac{-1\pm i\sqrt{3}}{2}.
\]
Substitution provides a direct check: each root satisfies
$x^2+x+1=0$.
\end{example}

The two roots in this example are complex conjugates. In general, the nonreal roots of a quadratic equation with real coefficients $a$, $b$, $c$
occur in conjugate pairs, while each real root coincides with its own conjugate.

Allowing complex solutions ensures that every quadratic equation has a solution. More generally, every polynomial equation of degree at least one,
\[
a_n x^n + a_{n-1} x^{n-1} + \cdots + a_1 x + a_0 = 0,
\]
has a complex solution. This result, proved by Gauss, is known as the \textit{Fundamental Theorem of Algebra}.

\section{Polar Representation and the Complex Exponential}

\subsection{Polar Coordinates: Modulus and Argument}

A complex number $z = a+bi$ corresponds to the point $(a,b)$ in the Argand plane. This point can also be described by polar coordinates
$(r,\theta)$, with $r \geq 0$, related to the Cartesian coordinates by
\[
\begin{aligned}
a &= r \cos \theta \\
b &= r \sin \theta.
\end{aligned}
\]
This gives the \emph{polar form} of $z$:
\[
\begin{aligned}
z &= a+bi \\
&= (r\cos\theta) + (r\sin\theta) i \\
&= r(\cos\theta + i \sin\theta),
\end{aligned}
\]
The modulus is given by
\[
\begin{aligned}
r &= |z| \\
&= (a^2+b^2)^{1/2}.
\end{aligned}
\]

When the real part is nonzero, the argument also satisfies $\tan \theta = b/a$.
For a nonzero complex number, the angle $\theta$ is called its \emph{argument}, written $\theta = \arg(z)$.
The argument is determined only up to an integer multiple of $2\pi$.

\begin{example}
We express $z = 1+i$ and $w = \sqrt{3} - i$ in polar form.
For $z$, we have
\[
\begin{gathered}
\begin{aligned}
r &= |z| \\
&= (1^2+1^2)^{1/2} \\
&= \sqrt{2}
\end{aligned} \\
\tan \theta = 1.
\end{gathered}
\]
Since $z$ lies in the first quadrant, we take $\theta = \pi/4$, giving
\[
z = \sqrt{2} ( \cos(\pi/4) + i \sin(\pi/4) ).
\]
For $w$, we have
\[
\begin{gathered}
\begin{aligned}
r &= |w| \\
&= (3+1)^{1/2} \\
&= 2
\end{aligned} \\
\tan \theta = -1/\sqrt{3}.
\end{gathered}
\]
Since $w$ lies in the fourth quadrant, we take $\theta = -\pi/6$, giving
\[
w = 2 ( \cos(-\pi/6) + i \sin(-\pi/6) ).
\]

\end{example}

Figure~\ref{fig:ch2-polar-form} shows the geometric meaning of modulus, argument, and Cartesian components.

\begin{figure}[!htbp]
\centering
\begin{tikzpicture}[scale=1.3, >=Stealth]
  \draw[->] (-0.5,0) -- (3.6,0) node[right] {$\mathrm{Re}$};
  \draw[->] (0,-0.5) -- (0,3.1) node[above] {$\mathrm{Im}$};
  \draw[figblue, thick, dashed] (0,0) circle (2.5cm);
  \node[figblue] at (2.9,2.15) {$|z|=r$};
  \coordinate (Z) at (50:2.5);
  \draw[vecB] (0,0) -- (Z) node[right, xshift=2pt] {$z = a+bi = r\,e^{i\theta}$};
  \draw[black!60, dashed] (Z) -- (Z |- 0,0) node[below] {$a = r\cos\theta$};
  \draw[black!60, dashed] (Z) -- (0,0 |- Z) node[left] {$b = r\sin\theta$};
  \draw[figgreen, thick] (0.6,0) arc (0:50:0.6);
  \node[figgreen] at (0.88,0.4) {$\theta$};
  \node[circle,fill=figred,inner sep=1.3pt] at (Z) {};
\end{tikzpicture}
\caption[Polar representation of a complex number]{Polar representation of
  $z = a+bi = r e^{i\theta} \in \C$ in the Argand--Gauss plane: modulus $r=|z|$,
  argument $\theta$, and Cartesian projections $a=r\cos\theta$, $b=r\sin\theta$.}
\label{fig:ch2-polar-form}
\end{figure}

\subsection{Arithmetic, Powers, and Roots in Polar Form}

Polar form reveals the geometry of multiplication and division. Given two complex numbers in polar form,
\[
\begin{aligned}
z_1 &= r_1(\cos\theta_1 + i\sin\theta_1) \\
z_2 &= r_2(\cos\theta_2 + i\sin\theta_2).
\end{aligned}
\]
Expanding their product and applying the addition formulas for sine and cosine yields $z_1 z_2 = r_1 r_2 \big[ \cos(\theta_1+\theta_2) + i \sin(\theta_1+\theta_2) \big]$, so multiplication multiplies the moduli and adds the arguments. Similarly, the subtraction formulas for sine and cosine show that division
divides the moduli and subtracts the arguments. For $z_2 \neq 0$, the quotient is
\[
\frac{z_1}{z_2} = \frac{r_1}{r_2} \big[ \cos(\theta_1-\theta_2) + i \sin(\theta_1-\theta_2) \big].
\]

Taking $z_1 = 1$, so that $r_1 = 1$ and $\theta_1 = 0$, and writing $z_2 = z$, gives the reciprocal formula. Any nonzero complex number in the following polar form has the reciprocal shown alongside it:
\[
\begin{aligned}
z &= r(\cos\theta + i \sin \theta), \\
\frac{1}{z} &= \frac{1}{r} \big( \cos\theta - i \sin\theta \big).
\end{aligned}
\]

\begin{example}
Using the polar representations of $1+i$ and $\sqrt{3}-i$ found above, the multiplication formula gives
\[
\begin{aligned}
(1+i)(\sqrt{3}-i) &= 2\sqrt{2} \left[ \cos(\pi/4 - \pi/6) + i \sin(\pi/4-\pi/6) \right] \\
&= 2\sqrt{2} \left[ \cos(\pi/12) + i \sin(\pi/12) \right].
\end{aligned}
\]
To verify the result in Cartesian form, multiply directly:
\[
\begin{aligned}
(1+i)(\sqrt{3}-i)
&=\sqrt{3}-i+i\sqrt{3}-i^2\\
&=(1+\sqrt{3})+i(\sqrt{3}-1).
\end{aligned}
\]
The two forms agree because
\[
2\sqrt{2}\cos(\pi/12)=1+\sqrt{3},\qquad
2\sqrt{2}\sin(\pi/12)=\sqrt{3}-1.
\]

\end{example}

Repeated multiplication in polar form gives a rule for powers. Consider a complex number in polar form:
\[
z = r(\cos\theta + i\sin\theta).
\]
Its square and cube are given by the following identities:
\[
\begin{aligned}
z^2 &= r^2(\cos 2\theta + i \sin 2\theta) \\
z^3 &= z^2 z \\
&= r^3(\cos 3\theta + i \sin 3\theta).
\end{aligned}
\]
The general result is named after the French mathematician Abraham de Moivre.

\begin{theorem}[De Moivre's Theorem]
Let $n$ be a positive integer. For a complex number in polar form, its $n$-th power is given by
\[
\begin{aligned}
z &= r(\cos\theta + i\sin\theta), \\
z^n &= \big[ r (\cos\theta + i \sin \theta) \big]^n \\
&= r^n \big( \cos n\theta + i \sin n\theta \big).
\end{aligned}
\]

\end{theorem}

Thus, taking the $n$-th power raises the modulus to the $n$-th power and multiplies the argument by $n$.

\begin{example}
We compute
\[
\left( \frac{1}{2} + \frac{1}{2}i \right)^{10}.
\]
The calculation uses the following factorization and the polar form of $1+i$ found above:
\[
\begin{aligned}
\frac{1}{2} + \frac{1}{2}i &= \frac{1}{2}(1+i), \\
1+i &= \sqrt{2}\big(\cos(\pi/4) + i\sin(\pi/4)\big).
\end{aligned}
\]
De Moivre's Theorem gives the tenth power from this modulus and argument.
Indeed,
\[
\frac12+\frac12i
=\frac{1}{\sqrt2}\left(\cos\frac{\pi}{4}+i\sin\frac{\pi}{4}\right),
\]
and therefore
\[
\begin{aligned}
\left(\frac12+\frac12i\right)^{10}
&=\left(\frac{1}{\sqrt2}\right)^{10}
\left(\cos\frac{10\pi}{4}+i\sin\frac{10\pi}{4}\right)\\
&=\frac1{32}\left(\cos\frac{5\pi}{2}+i\sin\frac{5\pi}{2}\right)\\
&=\frac{i}{32}.
\end{aligned}
\]
\end{example}

De Moivre's Theorem also gives the $n$-th roots of a complex number. An $n$-th root of $z$ is a complex number $w$ satisfying $w^n = z$.
For nonzero $z$, write both numbers in polar form:
\[
\begin{aligned}
w &= s(\cos\varphi + i\sin\varphi) \\
z &= r(\cos\theta + i\sin\theta).
\end{aligned}
\]
Applying De Moivre's Theorem to $w^n$, the equation $w^n = z$ becomes
\[
s^n(\cos n\varphi + i \sin n\varphi) = r(\cos\theta + i \sin\theta).
\]

Comparing moduli and arguments on both sides, and recalling that sine and cosine have period $2\pi$, gives the following relations for some integer $k$:
\[
\begin{aligned}
s &= r^{1/n} \\
n\varphi &= \theta + 2k\pi.
\end{aligned}
\]

This gives $\varphi = (\theta + 2k\pi)/n$. Each index $k = 0, 1, 2, \dots, n-1$ yields a distinct value of $w$, giving the following description of the $n$-th roots of a complex number.

\begin{theorem}[Roots of a Complex Number]
Let $n$ be a positive integer and let $z$ be a nonzero complex number with the polar representation below. For each index $k = 0, 1, \dots, n-1$, the second formula gives one of its exactly $n$ distinct $n$-th roots:
\[
\begin{aligned}
z &= r(\cos\theta + i\sin\theta), \\
w_k &= r^{1/n} \left[ \cos\left( \frac{\theta + 2k\pi}{n} \right) + i \sin\left( \frac{\theta+2k\pi}{n} \right) \right].
\end{aligned}
\]

\end{theorem}

All these roots have modulus $r^{1/n}$ and therefore lie on a circle of radius $r^{1/n}$ in the complex plane.
Their successive arguments differ by $2\pi/n$, so the roots are equally spaced around the circle.

\begin{example}
We find the sixth roots of $z = -8$. In trigonometric form,
\[
z = 8(\cos\pi + i\sin\pi).
\]
Applying the formula with $n = 6$, we obtain the six roots by taking $k = 0, 1, 2, 3, 4, 5$ in
\[
8^{1/6} \left[ \cos\left( \frac{\pi + 2k\pi}{6} \right) + i \sin\left( \frac{\pi+2k\pi}{6} \right) \right].
\]
All six resulting points lie on the circle of radius $8^{1/6} = \sqrt{2}$ in the complex plane, equally spaced around it.
Their arguments and Cartesian forms are
\[
\begin{aligned}
w_0&=\sqrt2\,e^{i\pi/6}
=\frac{\sqrt6}{2}+i\frac{\sqrt2}{2},\\
w_1&=\sqrt2\,e^{i\pi/2}=i\sqrt2,\\
w_2&=\sqrt2\,e^{i5\pi/6}
=-\frac{\sqrt6}{2}+i\frac{\sqrt2}{2},\\
w_3&=\sqrt2\,e^{i7\pi/6}
=-\frac{\sqrt6}{2}-i\frac{\sqrt2}{2},\\
w_4&=\sqrt2\,e^{i3\pi/2}=-i\sqrt2,\\
w_5&=\sqrt2\,e^{i11\pi/6}
=\frac{\sqrt6}{2}-i\frac{\sqrt2}{2}.
\end{aligned}
\]
For every $k$, raising $w_k$ to the sixth power multiplies its argument
by $6$ and raises its modulus $\sqrt{2}$ to the sixth power, yielding
$(\sqrt{2})^6=8$ and hence $w_k^6=-8$.
\end{example}

Figure~\ref{fig:sixth-roots-minus-eight} displays these six roots and the equal angular spacing between successive ones.

\begin{figure}[!htbp]
\centering
\begin{tikzpicture}[scale=2.1,>=Stealth]
  \draw[->,black!55] (-1.85,0) -- (1.95,0) node[right] {$\mathrm{Re}$};
  \draw[->,black!55] (0,-1.85) -- (0,1.95) node[above] {$\mathrm{Im}$};
  \draw[figblue,thick] (0,0) circle ({sqrt(2)});
  \draw[figblue!45,dashed] (30:{sqrt(2)}) -- (90:{sqrt(2)}) -- (150:{sqrt(2)})
    -- (210:{sqrt(2)}) -- (270:{sqrt(2)}) -- (330:{sqrt(2)}) -- cycle;
  \foreach \angle/\rootindex in {30/0,90/1,150/2,210/3,270/4,330/5} {
    \draw[gridline,dashed] (0,0) -- (\angle:{sqrt(2)});
    \fill[figblue] (\angle:{sqrt(2)}) circle (1.3pt);
  }
  \node[above right] at (30:{sqrt(2)}) {$w_0$};
  \node[above right] at (90:{sqrt(2)}) {$w_1$};
  \node[above left] at (150:{sqrt(2)}) {$w_2$};
  \node[below left] at (210:{sqrt(2)}) {$w_3$};
  \node[below right] at (270:{sqrt(2)}) {$w_4$};
  \node[below right] at (330:{sqrt(2)}) {$w_5$};
  \draw[figgreen,thick,->] (30:0.55) arc[start angle=30,end angle=90,radius=0.55];
  \node[figgreen] at (60:0.85) {$\pi/3$};
  \node[below left] at (0,0) {$0$};
\end{tikzpicture}
\caption[The six sixth roots of minus eight]{The solutions of $w^6=-8$ form the vertices of a regular hexagon on the circle of radius $\sqrt{2}$. Starting at argument $\pi/6$, the labels follow increasing argument in steps of $\pi/3$. Dashed segments emphasize the symmetry and do not represent additional roots.}
\label{fig:sixth-roots-minus-eight}
\end{figure}

\subsection{The Complex Exponential and Euler's Formula}


We now define $e^z$ for a complex number $z = x+iy$. Power series, studied in detail in later chapters, also allow complex terms.
Motivated by the Maclaurin series for $e^x$, we define the complex exponential by
\[
\begin{aligned}
e^z &:= \sum_{n=0}^{\infty} \frac{z^n}{n!} \\
&= 1 + z + \frac{z^2}{2!} + \frac{z^3}{3!} + \cdots.
\end{aligned}
\]
This has the same fundamental algebraic properties as the real exponential. In particular, $e^{z+w} = e^z e^w$, for all complex numbers $z$ and $w$.


Substituting $z = iy$ for real $y$ into the defining series and using $i^2 = -1$, $i^3 = -i$, $i^4 = 1$, together with the periodic recurrence of higher powers of $i$, one obtains
\[
\begin{aligned}
e^{iy} &= 1 + iy + \frac{(iy)^2}{2!} + \frac{(iy)^3}{3!} + \frac{(iy)^4}{4!} + \cdots \\
&= \left(1 - \frac{y^2}{2!} + \frac{y^4}{4!} - \cdots\right) + i \left( y - \frac{y^3}{3!} + \frac{y^5}{5!} - \cdots \right).
\end{aligned}
\]

The two series in parentheses are the Maclaurin series for $\cos y$ and $\sin y$, respectively. This gives Euler's formula.

\begin{theorem}[Euler's Formula]
For every real number $y$,
\[
e^{iy} = \cos y + i \sin y.
\]

\end{theorem}

Combining Euler's formula with the multiplicative property gives
\[
\begin{aligned}
e^{x+iy} &= e^x e^{iy} \\
&= e^x (\cos y + i \sin y).
\end{aligned}
\]
This expresses the exponential of $x+iy$ using the real exponential and trigonometric functions.

\begin{example}
We evaluate two particular complex exponentials using Euler's formula. First,
\[
\begin{aligned}
e^{i\pi} &= \cos\pi + i \sin\pi \\
&= -1 + i \cdot 0 \\
&= -1.
\end{aligned}
\]
This is an identity linking five fundamental mathematical constants. Second, the general formula for $e^{x+iy}$ gives
\[
\begin{aligned}
e^{1+i\pi/2} &= e^1 (\cos(\pi/2) + i \sin(\pi/2)) \\
&= e(0 + i \cdot 1) \\
&= ie.
\end{aligned}
\]

\end{example}

\enlargethispage{3\baselineskip}
Euler's formula also gives a short proof of De Moivre's Theorem. Writing $z = re^{i\theta}$, we obtain
\[
\begin{aligned}
\big[ r(\cos\theta + i\sin\theta) \big]^n &= \big( r e^{i\theta} \big)^n \\
&= r^n e^{in\theta} \\
&= r^n (\cos n\theta + i \sin n\theta),
\end{aligned}
\]
so De Moivre's Theorem follows directly from the laws of exponents for the complex exponential.
